\paragraph{6. Proof of Theorem 8.}
We can assume that $q\leq w$, for else the statistical bound is trivial and the theorem is true. We apply Theorem 14 with $\alpha=1/10$ and the sequence
\[t_i:=m/w^{c_0(q/\alpha)^{-c_1i}},\]
for constants $c_0,c_1$ to be set later. For large enough $c_1$ this satisfies the hypothesis of the theorem that $t_i\geq t_{i-1}\cdot cqw/\alpha$. We also need to show that $2^{-t_0}\geq\sqrt{8/|S|}$, where $|S|$ is the number of matrices $\Pi$ in the definition of $\operatorname{CMPred}$. This is true since $|S|\geq2^{\Omega(m/w)}$.
Let $k,D_k$, and $t_k$ be as provided by the theorem. Recall that
\[H_\infty(f(D_k))\geq H(\Pi)-t_{k-1}\cdot O(qw/\alpha)^2.\]
\paragraph{Finding comets among trees.}
The theorem provides $t:=t_k=m/w^{c_0(q/\alpha)^{-c_1k}}$ trees. Applying the Comet-Finding Lemma 21 with $c=3$ and $\ell=w\geq\log m$ gives a set of
\[t':=m/w^{c_0(q/\alpha)^{-c_1k}+O(1)}\]
disjoint $w^3$-comets, where the head lengths are in $[w^h,w^{h+1})$ for some $h\leq c_0(q/\alpha)+O(1)$. Note that to apply the lemma we need that $c_0(q/\alpha)\leq w$. This is guaranteed since $w\geq\log m$ and $q=O(\log m)/\log\log m$ for else the conclusion of the theorem holds trivially.
We shall get a contradiction looking at the row of the matrix corresponding to blocks of length $w^{h+2}$; the other rows can be ignored.
\paragraph{A random comet.}
To each of the above $t'$ comets we associate three relevant, consecutive blocks. Of these, the middle block is the first block that intersects the head of the comet. Note that:
\begin{itemize}
\item the relevant blocks cover the head of the comet, since the blocks have length $w^{h+2}$ while the head has length $\leq w^{h+1}$.
\item the relevant blocks of different comets are disjoint, since the tails of each comet have length $\geq w^h\cdot w^3$, while the blocks relevant to a comet are contained in an interval of length $3w^{h+2}$ intersecting the head.
\end{itemize}
Note that from $\operatorname{CMPred}(\Pi)$ we can reconstruct $\Pi$, and moreover $f(D_k)$ is in the range of $\operatorname{CMPred}$. Hence we can define
\[X:=\operatorname{CMPred}^{-1}(f(D_k))\]
and we have $H(X)=H(f(D_k))$. Let $B_i$ be the portion of $X$ in the three blocks relevant to comet $i$, in our current set of $t'$ comets. Recall that in row $h+2$ of the $\operatorname{CMPred}$ distribution $\Pi$, each block is given by a variable uniform over a support of size $2\cdot w^{h+2}$. Hence $B_i$ is a random variable uniform over its support $\operatorname{Supp}(B_i)$ of size $(2\cdot w^{h+2})^3$.
We want to argue that one such variable is close to uniform in our distribution $f(D_k)$. Indeed, recall from the beginning of the proof that
\[H(X)\geq H_\infty(f(D_k))\geq H(\Pi)-t_{k-1}\cdot O(qw/\alpha)^2.\]
Since $H(X,Y)\leq H(X)+H(Y)$ for any random variables $X,Y$, we have that,
\[H(B_1,B_2,\ldots,B_{t'})\geq t'\log|\operatorname{Supp}(B_i)|-t_{k-1}\cdot O(qw/\alpha)^2.\]
By Lemma 22, each $B_i$ is $\alpha$-close to uniform, except for those in a ``forbidden'' set of size $t_{k-1}\cdot O(q^2w^2/\alpha^4)$.
Now for the critical point, $t'$ is larger than the size of this forbidden set. This is true because we only lost $w^{O(1)}$ factors, so it suffices to make the constant $c_1$ large enough in the definition of the sequence $t_i$. Formally,
\[t_{k-1}\cdot O(q^2w^2/\alpha^4)=(m/w^{c_0(q/\alpha)^{-c_1(k-1)}})\cdot O(q^2w^2/\alpha^4)\]
which is smaller than $t'=m/w^{c_0(q/\alpha)^{-c_1k}+O(1)}$ for $c_1$ large enough. Here we are using that $q\leq O(\log m)/\log\log m$, $w\geq\log m$, $\alpha=\Theta(1)$.
\paragraph{Breaking correlation in the random comet.}
At this point we have a $w^3$-comet $(p,i,j)$ where the head length $(j-i)$ is in $[w^h,w^{h+1}]$ and
\begin{enumerate}
\item The answers to queries $i$ and $j$ are $\alpha$-independent, and
\item the relevant blocks are $\alpha$-close to uniform.
\end{enumerate}
In the query answers consider just the color corresponding to row $h+2$ for query $i$ and $j$. Let them be $C(i)$ and $C(j)$.
Because the relevant blocks are $\alpha$-close to uniform, for any color $c$ we have both $\mathbb P[C(i)=c]\leq1/2+\alpha$ and $\mathbb P[C(j)=c]\leq1/2+\alpha$. Also, because $C(i)$ and $C(j)$ are $\alpha$-independent, we have $\mathbb P[C(i)=C(j)]\leq1/2+2\alpha$.
However, $C(i)$ and $C(j)$ are in fact highly correlated. The only event in which $\mathbb P[C(i)\neq C(j)]$ is if the head of the comet contains an element. The head has length $\leq w^{h+1}$. The blocks have length $w^{h+2}$. If the variables in the blocks were uniform, the chance that the head contains an element is $\leq1/w$. The block is only $\alpha$-close to uniform, so this probability is $\leq1/w+\alpha$. Hence, $\mathbb P[C(i)=C(j)]\geq1-1/w-\alpha$. For $\alpha=1/10$, this is larger than the above value of $1/2-2\alpha$, concluding the proof.
